\centerline{\bf \Huge Серьезная олимпиада 1}

\begin{flushright}
        \scriptsize{Эпизод 3. Молчащий телефон}
\end{flushright}

\problem{1}
Существуют ли нечетные числа $k, l$ и $m$, такие что $k l+1, l m+1$ и $k m+1$ - полные квадраты?

\problem{2}
Окружности $\omega_1$ и $\omega_2$ с центрами $O_1$ и $O_2$ пересекаются в точках $A$ и $B$. На окружности $\omega_1$ выбрана точка $M$. Прямые $M A$ и $M B$ вторично пересекают окружность $\omega_2$ в точках $P$ и $Q$ соответственно. Докажите, что если четырехугольник $A O_1 B O_2$ вписанный, то точка пересечения прямых $A Q$ и $B P$ лежит на $\omega_1$.

\problem{3}
Петя и Вася играют в игру. Изначально дан пустой граф с $n$ вершинами. За один ход разрешается провести одно ориентированное ребро. Запрещается проводить кратные ребра (в том числе разнонаправленные), и запрещается делать граф сильно связным. Проигрывает тот, кто не может делать ход. Кто выигрывает при правильной игре?

\problem{4}
Для $n \in \mathbb{N}$ и положительных $a, b, c \in \mathbb{R}$ докажите неравенство:

$$
\frac{a^{n+1}}{a^n+a^{n-1} b+\ldots+b^n}+\frac{b^{n+1}}{b^n+b^{n-1} c+\ldots+c^n}+\frac{c^{n+1}}{c^n+c^{n-1} a+\ldots+a^n} \geqslant \frac{a+b+c}{n+1} .
$$


\problem{5}
Докажите, что для любого $n \in \mathbb{N}$ существует такое множество $X_n$, что $X_n \subset \mathbb{N}$, $\left|X_n\right|=n$ и для любого $M \subset X_n, M=\left\{m_1, \ldots, m_k\right\}$, выполнено

$$
\frac{m_1+\ldots+m_k}{k}=a^b
$$


для некоторых $a \in \mathbb{N}, b \in \mathbb{N}$ и $b>1$.